% =============================================================================
% Assignments/CS301/07-linked-lists-doubly-applications-coding-assignment.tex
% Coding Assignment 7: Looking Both Ways (P7, doubly-linked-list operations,
% plus the week's sparse-polynomial addition)
%
% Student-facing problem statement. Instructor-only files live in
% Assignments/CS301/07-linked-lists-doubly-applications-coding/:
%   reference_solution.c  (NEVER published)
%   test_hidden.c         (NEVER published)
%   autograder.sh         (NEVER published)
%   dlist.h, starter_dlist.c, test_visible.c
%
% Fresh instances throughout: no worked example re-asks a deck/notes,
% written-assignment, or earlier-lab instance verbatim (deck/notes: build
% 10-20-30; insert 20 after 10; delete 15 from 5-15-25; delete 5/25 from
% 5-15-25; poly (5x^3+4x+2)+(3x^2+4x+1); poly (4x^4+2x^2+7)+(3x^3+2x^2+5);
% poly (2x^2+1)+(-2x^2+5x)=5x+1; linked-stack 30-20-10; linked-queue
% 10-20-30). Visible tests use 7/14/21/28/42, 28/35 and
% (5x^2+3)+(-5x^2+4x)=4x+3 — all new.
%
% Compile with: cd Assignments/CS301 &&
%   TEXINPUTS=../../Preambles;$TEXINPUTS xelatex -interaction=nonstopmode 07-linked-lists-doubly-applications-coding-assignment.tex
%   (Windows/MiKTeX uses ; not : as the path separator)
% =============================================================================

\documentclass[11pt]{article}
\usepackage[margin=1in]{geometry}
% =============================================================================
% Preambles/header.tex — shared LaTeX/Beamer preamble for this project
%
% Use from a Beamer lecture by prepending:
%     \input{header}
% and compiling with TEXINPUTS=../Preambles:$TEXINPUTS (the /compile-latex
% skill does this automatically).
%
% PALETTE CONTRACT. Color names defined here MUST match the names in
% Quarto/theme-template.scss so Beamer and Quarto renderings use the same
% palette. The scripts/check-palette-sync.sh script enforces this.
%
% To customize: edit the HEX values below AND the matching $variable in
% Quarto/theme-template.scss, then run scripts/check-palette-sync.sh.
% =============================================================================

% -----------------------------------------------------------------------------
% Engine detection + encoding
%
% Our /compile-latex skill uses XeLaTeX, where inputenc is unnecessary and
% can produce encoding warnings. Only load inputenc under pdfTeX.
% -----------------------------------------------------------------------------
\usepackage{iftex}
\ifPDFTeX
  \usepackage[utf8]{inputenc}
\fi

\usepackage{amsmath, amssymb, amsthm}
\usepackage{booktabs}
\usepackage{graphicx}

% -----------------------------------------------------------------------------
% PALETTE — mirrors Quarto/theme-template.scss variable names.
% Load xcolor and define colors BEFORE anything that references them
% (notably hyperref, hypersetup, and the Beamer color assignments).
% -----------------------------------------------------------------------------
\usepackage{xcolor}

\definecolor{primary-blue}{HTML}{012169}      % headings, accents
\definecolor{primary-gold}{HTML}{B9975B}      % emphasis, borders
\definecolor{highlight-yellow}{HTML}{F2A900}  % markers, alerts
\definecolor{light-bg}{HTML}{E8EDF5}          % title / section backgrounds
\definecolor{jet}{HTML}{1A1A1A}               % body text

% Semantic colors (used in diagrams and annotations)
\definecolor{positive}{HTML}{15803D}          % good / observed / identified
\definecolor{negative}{HTML}{B91C1C}          % bad / problematic / bias
\definecolor{neutral}{HTML}{525252}           % reference / context

% Highlight accents (match .hi-* classes in the SCSS)
\definecolor{hi-slate}{HTML}{314F4F}
\definecolor{hi-green}{HTML}{2E8B57}
\definecolor{hi-red}{HTML}{C41E3A}

% Semantic aliases so skill templates and snippets can write
% \textcolor{accent}{...} without hard-coding "primary-blue".
\colorlet{accent}{primary-blue}
\colorlet{accent2}{primary-gold}

% -----------------------------------------------------------------------------
% Hyperref — loaded AFTER xcolor + palette so link colors resolve.
% -----------------------------------------------------------------------------
\usepackage{hyperref}
\hypersetup{colorlinks=true, linkcolor=primary-blue, urlcolor=primary-blue,
            citecolor=primary-blue, pdfborder={0 0 0}}

% -----------------------------------------------------------------------------
% Course code — set once per deck (after \input{header}, before \begin{document})
% via \coursecode{CS401}. Shown in the footer on every slide so a deck's course
% is identifiable without opening the title page. Empty by default.
% -----------------------------------------------------------------------------
\makeatletter
\newcommand{\coursecode}[1]{\gdef\@coursecodetext{#1}}
\gdef\@coursecodetext{}
\makeatother

% -----------------------------------------------------------------------------
% Beamer theme assignments (only apply under Beamer).
%
% We detect Beamer via \@ifclassloaded{beamer}, which is the documented,
% non-internal way. This must be wrapped in \makeatletter/\makeatother
% because \@ifclassloaded contains an @-catcode character.
% -----------------------------------------------------------------------------
\makeatletter
\@ifclassloaded{beamer}{%
  \usefonttheme{professionalfonts}
  \setbeamercolor{structure}{fg=primary-blue}
  \setbeamercolor{frametitle}{fg=primary-blue}
  \setbeamercolor{title}{fg=primary-blue}
  \setbeamercolor{subtitle}{fg=primary-gold}
  \setbeamercolor{author}{fg=primary-blue}
  \setbeamercolor{institute}{fg=neutral}
  \setbeamercolor{date}{fg=primary-gold}
  \setbeamercolor{itemize item}{fg=primary-blue}
  \setbeamercolor{itemize subitem}{fg=primary-gold}
  \setbeamercolor{alerted text}{fg=primary-gold}
  \setbeamercolor{block title}{fg=white, bg=primary-blue}
  \setbeamercolor{block body}{bg=light-bg}
  % Semantic block variants: block=definition/concept (blue), exampleblock=
  % worked example (green/positive), alertblock=key takeaway or caution (gold).
  % Keep to <=2 colored boxes per slide across ALL three variants (INV-7).
  \setbeamercolor{block title example}{fg=white, bg=positive}
  \setbeamercolor{block body example}{bg=light-bg}
  \setbeamercolor{block title alerted}{fg=white, bg=primary-gold}
  \setbeamercolor{block body alerted}{bg=light-bg}

  % Minimal footer: course code (left) + page X / N (right), no navigation clutter
  \setbeamertemplate{navigation symbols}{}
  \setbeamertemplate{footline}{%
    \color{neutral}\footnotesize\hspace{0.7em}\@coursecodetext\hfill
    \insertframenumber/\inserttotalframenumber\hspace{0.7em}\vspace{0.3em}}
}{}
\makeatother

% -----------------------------------------------------------------------------
% TikZ setup — libraries the snippet gallery uses
% -----------------------------------------------------------------------------
\usepackage{tikz}
\usetikzlibrary{arrows.meta, positioning, calc, decorations.pathreplacing,
                fit, shapes.geometric, backgrounds}

% Shared TikZ styles — snippets and hand-written diagrams can pick these up
% via `\begin{tikzpicture}[every node/.style={...}]` or by naming specific
% styles (e.g., `\node[dag-node] (X) at (0,0) {X};`).
\tikzset{
  % Node styles for causal DAGs and similar
  dag-node/.style = {
    draw=primary-blue, fill=light-bg, rounded corners=2pt,
    minimum width=1.2cm, minimum height=0.9cm, align=center, font=\small
  },
  decision-node/.style = {
    draw=primary-gold, fill=white, diamond, aspect=1.8,
    minimum width=1.8cm, minimum height=1.2cm, align=center, font=\small,
    inner sep=1pt
  },
  % Edge styles — observed vs counterfactual
  observed-edge/.style = {-{Latex[length=2.5mm]}, thick, draw=jet},
  counterfactual-edge/.style = {-{Latex[length=2.5mm]}, thick, draw=neutral, dashed},
  confound-edge/.style = {-{Latex[length=2.5mm]}, thick, draw=negative, dashed},
  % Data-point styles for plots
  observed-dot/.style = {fill=primary-blue, draw=primary-blue, circle, inner sep=1.2pt},
  counterfactual-dot/.style = {fill=white, draw=neutral, circle, inner sep=1.2pt}
}

% -----------------------------------------------------------------------------
% Convenience macros
% -----------------------------------------------------------------------------
\newcommand{\muted}[1]{\textcolor{neutral}{#1}}
\newcommand{\key}[1]{\textcolor{primary-gold}{\textbf{#1}}}
\newcommand{\good}[1]{\textcolor{positive}{#1}}
\newcommand{\bad}[1]{\textcolor{negative}{#1}}

% Standout frame macro: full-bleed dark background for section transitions.
% Beamer-only (uses \begin{frame}/\setbeamercolor/\usebeamerfont) -- do not
% call from an article-class file (e.g. Notes/<CODE>/*-notes.tex). \muted,
% \key, \good, \bad, and \coursecode above are plain xcolor/gdef and are
% safe to use from either class; verified when Notes/article-class support
% was added.
% Expand: \transitionslide{Your transition title}
\newcommand{\transitionslide}[1]{%
  \begingroup
  \setbeamercolor{background canvas}{bg=primary-blue}
  \begin{frame}[plain]
    \centering
    \vfill
    {\usebeamerfont{title}\color{white}\Large #1\par}
    \vfill
  \end{frame}
  \endgroup
}

% Section divider frame (Beamer-only), an alternative to \transitionslide with a
% darker two-line style: near-black (jet) background, a small white label line
% above a larger gold title, and a thin gold rule along the bottom edge.
% Expand: \sectiondivider{Section 2}{Pointers}
\newcommand{\sectiondivider}[2]{%
  \begingroup
  \setbeamercolor{background canvas}{bg=jet}
  \begin{frame}[plain]
    \vspace*{3.0cm}
    {\color{white}\ttfamily\large #1\par}
    \vspace{0.5cm}
    {\color{primary-gold}\LARGE #2\par}
    \begin{tikzpicture}[remember picture, overlay]
      \draw[primary-gold, line width=2.5pt]
        ($(current page.south west)+(0,0.1)$) -- ($(current page.south east)+(0,0.1)$);
    \end{tikzpicture}
  \end{frame}
  \endgroup
}

% -----------------------------------------------------------------------------
% Channel watermark — "Concepts That Click" (YouTube).
%
% Article-class files (CompetitiveExam Q/A sets, Notes, Assignments) get a
% light diagonal draft watermark on every page plus a centered footer naming
% the channel. Beamer decks get a subtle TikZ background watermark instead
% (the footline already carries the course code). Centralized here so every
% MCQ PDF is branded without per-file edits.
% -----------------------------------------------------------------------------
\makeatletter
\@ifclassloaded{article}{%
  \usepackage{draftwatermark}
  \SetWatermarkText{Concepts That Click}
  \SetWatermarkScale{1}
  \SetWatermarkFontSize{2cm}
  \SetWatermarkLightness{0.86}
  \SetWatermarkAngle{45}
  \usepackage{fancyhdr}
  \pagestyle{fancy}
  \fancyhf{}
  \fancyfoot[C]{\footnotesize\textcolor{neutral}{Concepts That Click~|~YouTube: \href{https://www.youtube.com/@concepts.thatclick}{@concepts.thatclick}}}
  \renewcommand{\headrulewidth}{0pt}
  \renewcommand{\footrulewidth}{0pt}
  \fancypagestyle{plain}{%
    \fancyhf{}%
    \fancyfoot[C]{\footnotesize\textcolor{neutral}{Concepts That Click~|~YouTube: \href{https://www.youtube.com/@concepts.thatclick}{@concepts.thatclick}}}%
    \renewcommand{\headrulewidth}{0pt}%
    \renewcommand{\footrulewidth}{0pt}%
  }
}{}
\@ifclassloaded{beamer}{%
  \setbeamertemplate{background}{%
    \begin{tikzpicture}[remember picture, overlay]
      \node[rotate=45, opacity=0.055, align=center]
        at (current page.center)
        {\Huge\textcolor{neutral}{Concepts That Click}};
    \end{tikzpicture}%
  }
}{}
\makeatother 
\coursecode{CS301}
\renewcommand{\thesection}{7.\arabic{section}}

\title{Coding Assignment 7: Looking Both Ways}
\author{Auqib Hamid Lone\\[0.3em]
  \emph{Department of Computer Science \& Engineering}, \emph{GCET Ganderbal}}
\date{\today}

\begin{document}
\maketitle

\noindent\textbf{Due:} two weeks from issue. There is no automatic late penalty --- if you need more time, ask before the deadline and an extension will be granted (see the course policy in the syllabus).

\section{Why this problem}

Week 6 built a singly linked list whose every delete paid a walk: to unlink a node you had to find the node whose \texttt{next} pointed at it, because the list could only look forward. Week 7 adds one pointer per node --- the predecessor link \texttt{prev} --- and that single change turns ``delete a node you already hold a pointer to'' from $O(n)$ into $O(1)$. This assignment (practical P7) is that doubly linked list, together with the week's headline application: sparse polynomials, where only the nonzero terms are stored and two polynomials are added by a single linear merge. You implement both on heap nodes, and every insert or delete must leave both links of the \emph{two-link invariant} in agreement: for every interior node \texttt{p}, \texttt{p->next->prev == p} and \texttt{p->prev->next == p}.

\section{Your task}

Implement the doubly linked list and sparse polynomial ADTs in C. The contract is given to you in \texttt{dlist.h} --- \emph{do not modify it}. Copy \texttt{starter\_dlist.c} to \texttt{dlist.c} and fill in all fifteen functions.

\textbf{Part A --- the doubly linked list.} The node is data plus a forward and a backward link:
\begin{verbatim}
typedef struct dnode { int data; struct dnode *prev; struct dnode *next; } DNode;

DNode *dlist_insert_front(DNode **head, int x);
DNode *dlist_insert_back(DNode **head, int x);
DNode *dlist_insert_after(DNode *p, int x);
int    dlist_delete_front(DNode **head, int *out);   /* 0 ok / -1 empty   */
int    dlist_delete_back(DNode **head, int *out);    /* 0 ok / -1 empty   */
int    dlist_delete_node(DNode **head, DNode *target);/* 0 ok / -1         */
DNode *dlist_search(DNode *head, int x);             /* NULL when absent  */
int    dlist_print_forward(const DNode *head, int *out, int out_cap);
int    dlist_print_backward(const DNode *head, int *out, int out_cap);
void   dlist_reverse(DNode **head);
void   dlist_free(DNode **head);
\end{verbatim}

Insertion is the deck's two-write pointer surgery, done for \emph{both} links. Inserting \texttt{x} after a known node \texttt{p} wires the new node between \texttt{p} and \texttt{p->next}, repairs the old successor's \texttt{prev} \emph{before} overwriting \texttt{p->next}, and is $O(1)$ worst case because the pointer \texttt{p} is already in hand; \texttt{dlist\_insert\_front} is the special case at \texttt{head}. Because this ADT caches no tail handle, \texttt{dlist\_insert\_back} and \texttt{dlist\_delete\_back} must walk to the last node, so each is $O(n)$ worst case --- state that honestly in your docstrings. \textbf{Deletion is the whole point of the backward link}: given only a pointer to a node, \texttt{dlist\_delete\_node} rewires \texttt{target->prev->next} and \texttt{target->next->prev} (moving \texttt{head} when the target is first), then frees the node exactly once, in $O(1)$ worst case. \texttt{print\_forward} and \texttt{print\_backward} copy the sequence into the caller's buffer (the ``print/display'' operation made testable, exactly as Week 5's \texttt{cqueue\_display} and Week 6's \texttt{slist\_display}); \texttt{print\_backward} walks to the tail and then back through \texttt{prev}. \texttt{reverse} flips \emph{both} links of every node.

\textbf{Part B --- sparse polynomials.} A polynomial $\sum_i c_i x^{e_i}$ is stored as a chain of \texttt{pnode} terms in strictly descending exponent order; a zero coefficient is never stored (the sparsity payoff). The polynomial chain is \emph{singly} linked with a cached \texttt{tail} --- document why: the merge only ever moves forward, so a \texttt{prev} link would buy nothing here, while the cached tail makes every append $O(1)$.
\begin{verbatim}
struct pnode { int coeff; int exp; struct pnode *next; };
struct poly  { struct pnode *head; struct pnode *tail; };

void poly_init(struct poly *p);
int  poly_append_term(struct poly *p, int coeff, int exp);  /* 0 ok / -1 fail */
void poly_free(struct poly *p);
int  poly_add(const struct poly *a, const struct poly *b, struct poly *out);
\end{verbatim}

\texttt{poly\_add} is a single linear merge of the two descending-exponent chains: copy the larger exponent through, add the coefficients on an equal exponent, suppress a term whose coefficient sums to zero, and when one chain is exhausted append the rest of the other. Build the output in descending order with $O(1)$ tail appends, so the whole addition is $O(n+m)$ worst case for $n$ and $m$ terms. An all-cancelling sum is not an error --- it leaves \texttt{out} as the zero polynomial (\texttt{head == NULL}). \texttt{out} must be the zero polynomial on entry and must not alias \texttt{a} or \texttt{b}; on any failure free partial work and leave \texttt{out} empty.

\section{Constraints}

\begin{itemize}
  \item Representation must be the heap-chained \texttt{DNode} (Part A) and \texttt{pnode} (Part B) nodes of \texttt{dlist.h} --- no arrays, no global/static mutable state.
  \item Every allocation is checked (\texttt{malloc} for \texttt{NULL}); every node is freed exactly once; no use-after-free, no writes past allocated bounds. Error paths leave the list or polynomial unchanged.
  \item The two-link invariant must hold after every list operation; the hidden suite re-checks it after each one.
  \item Correctness is graded against the contract; your code must compile with \texttt{gcc -std=c11 -Wall -Wextra -Werror} without errors or warnings.
  \item Complexity: \texttt{insert\_front}, \texttt{insert\_after}, \texttt{delete\_front}, and \texttt{delete\_node} are $O(1)$ worst case; \texttt{insert\_back} and \texttt{delete\_back} are $O(n)$ worst case (no tail handle); \texttt{search} is $O(n)$ worst case and $O(1)$ best case; both prints are $\Theta(n)$; \texttt{reverse} is $O(n)$ worst case in $O(1)$ auxiliary space; \texttt{poly\_append\_term} is $O(1)$ worst case; \texttt{poly\_add} is $O(n+m)$ worst case.
\end{itemize}

\section{Worked examples (these are the visible tests)}

\textbf{Example 1 --- build the chain and read it both ways.} End-insert $14$, then $28$, then $42$. Front-insert $7$ (chain $7 \leftrightarrow 14 \leftrightarrow 28 \leftrightarrow 42$). Insert $21$ after the node holding $14$ (chain $7 \leftrightarrow 14 \leftrightarrow 21 \leftrightarrow 28 \leftrightarrow 42$). \texttt{print\_forward} writes $7, 14, 21, 28, 42$; \texttt{print\_backward} writes $42, 28, 21, 14, 7$. Freeing sets \texttt{head} to \texttt{NULL}.

\textbf{Example 2 --- delete from both ends and by pointer.} From $14 \leftrightarrow 28 \leftrightarrow 42$: \texttt{delete\_front} returns $14$ (chain $28 \leftrightarrow 42$), then \texttt{delete\_back} returns $42$ (chain $28$). \texttt{search(28)} finds the node; \texttt{search(99)} reports absent (\texttt{NULL}). Deleting that node by pointer empties the list; \texttt{delete\_front} on the empty list reports failure.

\textbf{Example 3 --- reverse, and a polynomial addition with cancellation.} Reversing $28 \leftrightarrow 35$ gives $35 \leftrightarrow 28$; \texttt{print\_backward} then reads $28, 35$. Reversing again restores $28 \leftrightarrow 35$. Adding $A = 5x^{2} + 3$ and $B = -5x^{2} + 4x$ gives $4x + 3$: the equal-exponent $x^{2}$ terms cancel to coefficient $0$ and no zero term is stored, while the unshared $4x$ and the constants copy through.

\section{What to submit}

\begin{enumerate}
  \item \texttt{dlist.c} --- your completed implementation (edit \texttt{starter\_dlist.c}).
  \item \texttt{dlist.h} --- unchanged, but include it so the submission compiles standalone.
\end{enumerate}

\section{Getting started}

\begin{enumerate}
  \item Copy \texttt{starter\_dlist.c} to \texttt{dlist.c} and fill in the fifteen function bodies.
  \item Sanity-check locally:
\begin{verbatim}
gcc -std=c11 -Wall -Wextra -Werror -o test_visible test_visible.c dlist.c
./test_visible
\end{verbatim}
    All three worked examples should pass.
  \sloppy
  \item Submit \texttt{dlist.c} and \texttt{dlist.h}. The autograder compiles your \texttt{dlist.c} against a \emph{hidden} test suite and reports pass/fail. The hidden suite covers: \texttt{NULL} head-pointer safety for every operation; front/back/after insert boundaries with neighbour-link repair and returned-node identity; delete of head/middle/tail by pointer using saved node pointers; \texttt{delete\_node(NULL)} and empty-list refusals with \texttt{*out} untouched; \texttt{out}-is-\texttt{NULL} calls; \texttt{print\_forward}/\texttt{print\_backward} on empty/single/many with a backward print that is the reverse of the forward print; buffer-cap and \texttt{NULL}-buffer refusals with a canary check that nothing is written, plus exact-capacity success and head purity; search identity at head and tail; reverse of empty/single/two/four chains with double-reverse identity; the two-link invariant re-verified after every operation; \texttt{INT\_MIN}/\texttt{INT\_MAX}/\texttt{-1}/\texttt{0} round-trips; polynomial helper boundaries (zero-coefficient skip, descending-order enforcement, double-free safety); \texttt{poly\_add} argument validation, empty operands, full and partial cancellation, all-zero results, different lengths, and \texttt{INT} extremes; and two seeded random campaigns checked operation-by-operation against independent oracles --- the list against an int-array model, the polynomial against an exponent-sum model.
\end{enumerate}

\end{document}
